\chapter{2007年辽宁卷（文）}

\section{单选题}

\begin{problem}
若集合 \(A = \{1, 3\}\)，\(N = \{2, 3, 4\}\)，则 \(A \cap B =\)（\quad）

\choices
    {\(\{1\}\)}
    {\(\{2\}\)}
    {\(\{3\}\)}
    {\(\{1, 2, 3, 4\}\)}
\end{problem}

\begin{answer}
C
\end{answer}

\begin{solution}
原卷对第二个集合的定义使用了符号 \(N\)，而交集结论处印作 \(A\cap B\). 按原卷给出的集合定义计算，有
\(A\cap N=\{1,3\}\cap\{2,3,4\}=\{3\}\). 因此按题面实际数据，选择 C.
\end{solution}

\begin{problem}
若函数 \(y = f(x)\) 的反函数图象过点 \((1, 5)\)，则函数 \(y = f(x)\) 的图象必过点（\quad）

\choices
    {\((5, 1)\)}
    {\((1, 5)\)}
    {\((1, 1)\)}
    {\((5, 5)\)}
\end{problem}

\begin{answer}
A
\end{answer}

\begin{solution}
函数与其反函数的图象关于直线 \(y=x\) 对称. 反函数图象上的点 \((1,5)\) 关于直线 \(y=x\) 的对称点为 \((5,1)\)，所以原函数图象必过 \((5,1)\)，选择 A.
\end{solution}

\begin{problem}
双曲线 \(\frac{x^2}{16} - \frac{y^2}{9} = 1\) 的焦点坐标为（\quad）

\choices
    {\((-\sqrt{7}, 0)\)，\((\sqrt{7}, 0)\)}
    {\((0, -\sqrt{7})\)，\((0, \sqrt{7})\)}
    {\((-5, 0)\)，\((5, 0)\)}
    {\((0, -5)\)，\((0, 5)\)}
\end{problem}

\begin{answer}
C
\end{answer}

\begin{solution}
由双曲线标准方程可知 \(a^2=16\)，\(b^2=9\). 双曲线的焦半距满足
\(c^2=a^2+b^2=16+9=25\)，所以 \(c=5\). 由于实轴在 \(x\) 轴上，两个焦点为 \((-5,0)\)，\((5,0)\)，选择 C.
\end{solution}

\begin{problem}
若向量 \(\bs{a}\) 与 \(\bs{b}\) 不共线，\(\bs{a}\cdot\bs{b}\neq0\)，且 \(\bs{c}=\bs{a}-\left(\frac{\bs{a}\cdot\bs{a}}{\bs{a}\cdot\bs{b}}\right)\bs{b}\)，则向量 \(\bs{a}\) 与 \(\bs{c}\) 的夹角为（\quad）

\choices
    {\(0\)}
    {\(\frac{\pi}{6}\)}
    {\(\frac{\pi}{3}\)}
    {\(\frac{\pi}{2}\)}
\end{problem}

\begin{answer}
D
\end{answer}

\begin{solution}
由题设，
\[
\bs{a}\cdot\bs{c}=\bs{a}\cdot\bs{a}-\frac{\bs{a}\cdot\bs{a}}{\bs{a}\cdot\bs{b}}\left(\bs{a}\cdot\bs{b}\right)=0
\]
又因为 \(\bs{a}\) 与 \(\bs{b}\) 不共线，所以 \(\bs{c}\neq\bs{0}\)，从而 \(\bs{a}\) 与 \(\bs{c}\) 垂直，夹角为 \(\frac{\pi}{2}\)，选择 D.
\end{solution}

\begin{problem}
设等差数列 \(\{a_n\}\) 的前 \(n\) 项和为 \(S_n\)，若 \(S_3=9\)，\(S_6=36\)，则 \(a_7+a_8+a_9=\)（\quad）

\choices
    {\(63\)}
    {\(45\)}
    {\(36\)}
    {\(27\)}
\end{problem}

\begin{answer}
B
\end{answer}

\begin{solution}
设等差数列首项为 \(a_1\)，公差为 \(d\). 由等差数列前 \(n\) 项和公式，
\[
S_3=3a_1+3d=9
\qquad
S_6=6a_1+15d=36
\]
于是 \(a_1+d=3\)，\(2a_1+5d=12\). 由 \(a_1=3-d\) 代入第二式，得
\(2(3-d)+5d=12\)，解得 \(d=2\)，从而 \(a_1=1\). 因此
\[
a_7+a_8+a_9=3a_8=3\left(a_1+7d\right)=3(1+14)=45
\]
所以选择 B.
\end{solution}

\begin{problem}
若 \(m\)，\(n\) 是两条不同的直线，\(\alpha\)，\(\beta\)，\(\gamma\) 是三个不同的平面，则下列命题中的真命题是（\quad）

\choices
    {若 \(m\subset\beta\)，\(\alpha\perp\beta\)，则 \(m\perp\alpha\)}
    {若 \(m\perp\beta\)，\(m\parallel\alpha\)，则 \(\alpha\perp\beta\)}
    {若 \(\alpha\perp\gamma\)，\(\alpha\perp\beta\)，则 \(\beta\perp\gamma\)}
    {若 \(\alpha\cap\gamma=m\)，\(\beta\cap\gamma=n\)，\(m\parallel n\)，则 \(\alpha\parallel\beta\)}
\end{problem}

\begin{answer}
B
\end{answer}

\begin{solution}
若直线 \(m\perp\beta\)，则 \(m\) 的方向是平面 \(\beta\) 的法向方向；又因为 \(m\parallel\alpha\)，所以平面 \(\alpha\) 内有一条直线平行于平面 \(\beta\) 的法线，即平面 \(\alpha\) 垂直于平面 \(\beta\). 因此 B 正确.

其余命题均不成立. 例如取 \(\beta:z=0\)，\(\alpha:y=0\)，令 \(m\) 为 \(x\) 轴，则 \(m\subset\beta\) 且 \(\alpha\perp\beta\)，但 \(m\subset\alpha\)，故 A 错误. 取 \(\alpha:z=0\)，\(\beta:y=0\)，\(\gamma:x+y=0\)，则 \(\alpha\perp\beta\) 且 \(\alpha\perp\gamma\)，而 \(\beta\) 与 \(\gamma\) 不垂直，故 C 错误. 取 \(\gamma:z=0\)，\(\alpha:x=0\)，\(\beta:x+z=1\)，则 \(\alpha\cap\gamma\) 与 \(\beta\cap\gamma\) 是两条平行直线，但 \(\alpha\) 与 \(\beta\) 不平行，故 D 错误. 所以选择 B.
\end{solution}

\begin{problem}
若函数 \(y=f(x)\) 的图象按向量 \(\bs{a}\) 平移后，得到函数 \(y=f(x-1)-2\) 的图象，则向量 \(\bs{a}=\)（\quad）

\choices
    {\((1,-2)\)}
    {\((1,2)\)}
    {\((-1,-2)\)}
    {\((-1,2)\)}
\end{problem}

\begin{answer}
A
\end{answer}

\begin{solution}
函数图象按向量 \((h,k)\) 平移后，方程由 \(y=f(x)\) 变为 \(y=f(x-h)+k\). 与题给方程
\(y=f(x-1)-2\) 比较，得 \(h=1\)，\(k=-2\)，所以 \(\bs{a}=(1,-2)\)，选择 A.
\end{solution}

\begin{problem}
已知变量 \(x\)，\(y\) 满足约束条件 \(\begin{cases}
x-y+2\leqslant0,\\
x\geqslant1,\\
x+y-7\leqslant0,
\end{cases}\) 则 \(\frac{y}{x}\) 的取值范围是（\quad）

\choices
    {\(\left[\frac{9}{5},6\right]\)}
    {\(\left(-\infty,\frac{9}{5}\right]\cup[6,+\infty)\)}
    {\((-\infty,3]\cup[6,+\infty)\)}
    {\([3,6]\)}
\end{problem}

\begin{answer}
A
\end{answer}

\begin{solution}
约束条件等价于
\(x\geqslant1\)，\(y\geqslant x+2\)，\(y\leqslant7-x\). 由上下界能够同时成立，得
\(x+2\leqslant7-x\)，所以 \(1\leqslant x\leqslant\frac{5}{2}\). 特别地，\(x>0\). 于是
\[
1+\frac{2}{x}\leqslant\frac{y}{x}\leqslant\frac{7}{x}-1
\]
左端随 \(x\) 增大而减小，右端也随 \(x\) 增大而减小，因此
\[
\frac{y}{x}\geqslant1+\frac{2}{5/2}=\frac{9}{5}
\qquad
\frac{y}{x}\leqslant\frac{7}{1}-1=6
\]
区域连通且函数 \(y/x\) 连续，故值域为 \(\left[\frac{9}{5},6\right]\)，选择 A.
\end{solution}

\begin{problem}
函数 \(\log_{\frac{1}{2}}(x^2-5x+6)\) 的单调减区间为（\quad）

\choices
    {\(\left(\frac{5}{2},+\infty\right)\)}
    {\((3,+\infty)\)}
    {\(\left(-\infty,\frac{5}{2}\right)\)}
    {\((-\infty,2)\)}
\end{problem}

\begin{answer}
B
\end{answer}

\begin{solution}
函数定义域由 \(x^2-5x+6=(x-2)(x-3)>0\) 给出，即 \(x<2\) 或 \(x>3\). 当 \(x>3\) 时，二次函数 \(x^2-5x+6\) 单调递增，而底数 \(\frac{1}{2}\) 的对数函数单调递减，所以复合函数在 \((3,+\infty)\) 上单调递减. 故选择 B.
\end{solution}

\begin{problem}
一个坛子里有编号为 \(1\)，\(2\)，\(\cdots\)，\(12\) 的 \(12\) 个大小相同的球，其中 \(1\) 到 \(6\) 号球是红球，其余的是黑球. 若从中任取两个球，则取到的都是红球，且至少有 \(1\) 个球的号码是偶数的概率为（\quad）

\choices
    {\(\frac{1}{22}\)}
    {\(\frac{1}{11}\)}
    {\(\frac{3}{22}\)}
    {\(\frac{2}{11}\)}
\end{problem}

\begin{answer}
D
\end{answer}

\begin{solution}
从 \(12\) 个球中任取两个球的基本事件数为 \(\mathrm C_{12}^{2}=66\). 满足两个球都是红球的选法有 \(\mathrm C_6^2\) 种，其中两个号码都是奇数的选法有 \(\mathrm C_3^2\) 种. 因此满足题意的选法数为
\[
\mathrm C_6^2-\mathrm C_3^2=15-3=12
\]
所求概率为 \(\frac{12}{66}=\frac{2}{11}\)，选择 D.
\end{solution}

\begin{problem}
设 \(p\)，\(q\) 是两个命题，\(p\)：\(|x|-3>0\)，\(q\)：\(x^2-\frac{5}{6}x+\frac{1}{6}>0\)，则 \(p\) 是 \(q\) 的（\quad）

\choices
    {充分而不必要条件}
    {必要而不充分条件}
    {充分必要条件}
    {既不充分也不必要条件}
\end{problem}

\begin{answer}
A
\end{answer}

\begin{solution}
由 \(|x|-3>0\) 得 \(x<-3\) 或 \(x>3\). 又
\[
x^2-\frac{5}{6}x+\frac{1}{6}=\frac{1}{6}(3x-1)(2x-1)
\]
所以 \(q\) 的解集为 \(x<\frac{1}{3}\) 或 \(x>\frac{1}{2}\). 显然 \(p\) 的解集包含于 \(q\) 的解集，故 \(p\) 是 \(q\) 的充分条件. 取 \(x=0\)，满足 \(q\) 但不满足 \(p\)，所以不是必要条件，选择 A.
\end{solution}

\begin{problem}
将数字 \(1\)，\(2\)，\(3\)，\(4\)，\(5\)，\(6\) 排成一列，记第 \(i\) 个数为 \(a_i\)（\(i=1,2,\cdots,6\)）. 若 \(a_1\neq1\)，\(a_3\neq3\)，\(a_5\neq5\)，\(a_1<a_3<a_5\)，则不同的排列方法种数为（\quad）

\choices
    {\(18\)}
    {\(30\)}
    {\(36\)}
    {\(48\)}
\end{problem}

\begin{answer}
B
\end{answer}

\begin{solution}
先选择放在第 \(1\)，\(3\)，\(5\) 个位置上的三个数，并按递增顺序放置. 三个数的集合共有 \(\mathrm C_6^3=20\) 种.

其中 \(a_1=1\) 的有 \(\mathrm C_5^2=10\) 种，\(a_3=3\) 的有 \(2\times\mathrm C_3^2=6\) 种，\(a_5=5\) 的有 \(\mathrm C_4^2=6\) 种. 两两同时发生的情形分别有 \(3\)，\(3\)，\(2\) 种，三者同时发生的情形有 \(1\) 种. 因此符合三个限制条件的集合数为
\[
20-(10+6+6)+(3+3+2)-1=5
\]
确定奇数位置的数后，剩余三个数在第 \(2\)，\(4\)，\(6\) 个位置上有 \(3!=6\) 种排法，所以总数为 \(5\times6=30\)，选择 B.
\end{solution}

\section{填空题}

\begin{problem}
已知函数 \(y=f(x)\) 为奇函数，若 \(f(3)-f(2)=1\)，则 \(f(-2)-f(-3)=\fillinblank{}\).
\end{problem}

\begin{answer}
\(1\)
\end{answer}

\begin{solution}
因为 \(f(x)\) 是奇函数，所以 \(f(-2)=-f(2)\)，\(f(-3)=-f(3)\). 因此
\[
f(-2)-f(-3)=-f(2)+f(3)=f(3)-f(2)=1
\]
\end{solution}

\begin{problem}
\(\left(\sqrt{x}+\frac{1}{\sqrt[4]{x}}\right)^8\) 展开式中含有 \(x\) 的整数次幂的项的系数之和为\fillinblank{}. （用数字作答）
\end{problem}

\begin{answer}
\(72\)
\end{answer}

\begin{solution}
设通项中取 \(k\) 个因子 \(\frac{1}{\sqrt[4]{x}}\)，则该项为
\[
\mathrm C_8^k\left(\sqrt{x}\right)^{8-k}\left(\frac{1}{\sqrt[4]{x}}\right)^k
=\mathrm C_8^k x^{4-\frac{3k}{4}}
\qquad k=0,1,\cdots,8
\]
指数 \(4-\frac{3k}{4}\) 为整数，当且仅当 \(k\) 是 \(4\) 的倍数. 因此 \(k=0,4,8\)，对应项的系数之和为
\[
\mathrm C_8^0+\mathrm C_8^4+\mathrm C_8^8=1+70+1=72
\]
\end{solution}

\begin{problem}
若一个底面边长为 \(\frac{\sqrt{6}}{2}\)，侧棱长为 \(\sqrt{6}\) 的正六棱柱的所有顶点都在一个球面上，则此球的体积为\fillinblank{}.
\end{problem}

\begin{answer}
\(4\pi\sqrt{3}\)
\end{answer}

\begin{solution}
正六边形的外接圆半径等于边长，所以底面外接圆半径为 \(R=\frac{\sqrt{6}}{2}\). 球心在棱柱两底面中心连线的中点，设球半径为 \(r\)，则
\[
r^2=R^2+\left(\frac{\sqrt{6}}{2}\right)^2
=\frac{6}{4}+\frac{6}{4}=3
\]
所以 \(r=\sqrt{3}\)，球的体积为
\[
V=\frac{4}{3}\pi r^3=\frac{4}{3}\pi\left(\sqrt{3}\right)^3=4\pi\sqrt{3}
\]
\end{solution}

\begin{problem}
设椭圆 \(\frac{x^2}{25}+\frac{y^2}{16}=1\) 上一点 \(P\) 到左准线的距离为 \(10\)，\(F\) 是该椭圆的左焦点. 若点 \(M\) 满足 \(\overrightarrow{OM}=\frac{1}{2}\left(\overrightarrow{OP}+\overrightarrow{OF}\right)\)，则 \(\left|\overrightarrow{OM}\right|=\fillinblank{}\).
\end{problem}

\begin{answer}
\(2\)
\end{answer}

\begin{solution}
椭圆的长半轴、短半轴分别为 \(a=5\)，\(b=4\)，故焦半距 \(c=\sqrt{a^2-b^2}=3\)，离心率 \(e=\frac{c}{a}=\frac{3}{5}\). 左准线方程为 \(x=-\frac{a}{e}=-\frac{25}{3}\)，左焦点为 \(F=(-3,0)\).

设 \(P=(x,y)\). 由于 \(P\) 在椭圆上，\(-5\leqslant x\leqslant5\)，从而 \(x+\frac{25}{3}>0\). 点 \(P\) 到左准线的距离为 \(10\)，所以
\[
x+\frac{25}{3}=10,\qquad x=\frac{5}{3}
\]
代入椭圆方程，得
\[
\frac{(5/3)^2}{25}+\frac{y^2}{16}=1,\qquad \frac{1}{9}+\frac{y^2}{16}=1,\qquad y^2=\frac{128}{9}
\]
由向量关系，
\[
\overrightarrow{OM}=\frac{1}{2}\left((x,y)+(-3,0)\right)
=\left(-\frac{2}{3},\frac{y}{2}\right)
\]
因此
\[
\left|\overrightarrow{OM}\right|^2=\frac{4}{9}+\frac{y^2}{4}=\frac{4}{9}+\frac{32}{9}=4
\]
所以 \(\left|\overrightarrow{OM}\right|=2\).
\end{solution}

\section{解答题}

\begin{problem}
某公司在过去几年内使用某种型号的灯管 \(1000\) 支，该公司对这些灯管的使用寿命（单位：小时）进行了统计，统计结果如下表所示：

\begin{center}\small
\begin{tabular}{|c|c|c|c|c|}
\hline
分组 & \([500,900)\) & \([900,1100)\) & \([1100,1300)\) & \([1300,1500)\) \\
\hline
频数 & 48 & 121 & 208 & 223 \\
\hline
频率 &  &  &  &  \\
\hline
\end{tabular}

\vspace{0.5em}

\begin{tabular}{|c|c|c|c|}
\hline
分组 & \([1500,1700)\) & \([1700,1900)\) & \([1900,+\infty)\) \\
\hline
频数 & 193 & 165 & 42 \\
\hline
频率 &  &  &  \\
\hline
\end{tabular}
\end{center}

\begin{enumerate}
\item 将各组的频率填入表中；
\item 根据上述统计结果，计算灯管使用寿命不足 \(1500\) 小时的频率；
\item 该公司某办公室新安装了这种型号的灯管 \(3\) 支，若将上述频率作为概率，试求至少有 \(2\) 支灯管的寿命不足 \(1500\) 小时的概率.
\end{enumerate}
\end{problem}

\begin{solution}
\begin{enumerate}
\item 各组频率等于相应频数除以总数 \(1000\)，所以从左到右依次为
\(0.048\)，\(0.121\)，\(0.208\)，\(0.223\)，\(0.193\)，\(0.165\)，\(0.042\). 这些数的和为 \(1\)，与总频率一致.
\item 使用寿命不足 \(1500\) 小时包括前四组，因此其频率为
\[
\frac{48+121+208+223}{1000}=\frac{600}{1000}=0.6
\]
\item 设 \(X\) 表示 \(3\) 支灯管中寿命不足 \(1500\) 小时的支数，则 \(X\) 服从参数为 \(3\)，\(0.6\) 的二项分布. 于是
\[
P(X\geqslant2)=\mathrm C_3^2(0.6)^2(0.4)+(0.6)^3
=0.432+0.216=0.648
\]
\end{enumerate}
\end{solution}

\begin{problem}
如图，在直三棱柱 \(ABC-A_1B_1C_1\) 中，\(\angle ACB=90^{\circ}\)，\(AC=BC=a\)，\(D\)，\(E\) 分别为棱 \(AB\)，\(BC\) 的中点，\(M\) 为棱 \(AA_1\) 上的点，二面角 \(M-DE-A\) 为 \(30^{\circ}\).

\begin{enumerate}
\item 证明：\(A_1B_1\perp C_1D\)；
\item 求 \(MA\) 的长，并求点 \(C\) 到平面 \(MDE\) 的距离.
\end{enumerate}

\bitmapfigure[max width=0.38\linewidth]{img/2007/liaoning_liberal/q18_fig1.png}
\end{problem}

\begin{solution}
建立空间直角坐标系，使 \(C=(0,0,0)\)，\(A=(a,0,0)\)，\(B=(0,a,0)\)，并设 \(AA_1=h\). 令 \(MA=m\)，则
\[
D=\left(\frac{a}{2},\frac{a}{2},0\right),
\quad E=\left(0,\frac{a}{2},0\right),
\]
\[
C_1=(0,0,h),
\quad A_1=(a,0,h),
\quad B_1=(0,a,h),
\quad M=(a,0,m)
\]

\begin{enumerate}
\item 有
\[
\overrightarrow{A_1B_1}=(-a,a,0),
\qquad
\overrightarrow{C_1D}=\left(\frac{a}{2},\frac{a}{2},-h\right)
\]
因此
\[
\overrightarrow{A_1B_1}\cdot\overrightarrow{C_1D}
=-\frac{a^2}{2}+\frac{a^2}{2}=0
\]
故 \(A_1B_1\perp C_1D\).
\item 直线 \(DE\) 平行于 \(x\) 轴. 平面 \(MDE\) 的方程可由点 \(D\)，\(E\)，\(M\) 求得，为
\[
2my+az-am=0
\]
底面平面 \(ADE\) 的方程为 \(z=0\). 在垂直于 \(DE\) 的截面内，二面角的正切为
\[
\tan30^{\circ}=\frac{m}{a/2}=\frac{2m}{a}
\]
所以 \(m=\frac{a}{2\sqrt{3}}=\frac{\sqrt{3}}{6}a\).

点 \(C=(0,0,0)\) 到平面 \(MDE\) 的距离为
\[
\frac{|{-am}|}{\sqrt{(2m)^2+a^2}}
=\frac{am}{\sqrt{4m^2+a^2}}
=\frac{a}{4}
\]
故 \(MA=\frac{a}{2\sqrt{3}}\)，点 \(C\) 到平面 \(MDE\) 的距离为 \(\frac{a}{4}\).
\end{enumerate}
\end{solution}

\begin{problem}
已知函数 \(f(x)=\sin\left(\omega x+\frac{\pi}{6}\right)+\sin\left(\omega x-\frac{\pi}{6}\right)-2\cos^2\frac{\omega x}{2}\)，\(x\in\R\)（其中 \(\omega>0\)）.

\begin{enumerate}
\item 求函数 \(f(x)\) 的值域；
\item 若函数 \(y=f(x)\) 的图象与直线 \(y=-1\) 的两个相邻交点间的距离为 \(\frac{\pi}{2}\)，求函数 \(y=f(x)\) 的单调区间.
\end{enumerate}
\end{problem}

\begin{solution}
令 \(t=\omega x\). 利用和差化积公式及降幂公式，
\[
f(x)=2\sin t\cos\frac{\pi}{6}-\left(1+\cos t\right)=\sqrt{3}\sin t-\cos t-1=2\sin\left(t-\frac{\pi}{6}\right)-1
\]
\begin{enumerate}
\item 因为 \(-1\leqslant\sin\left(t-\frac{\pi}{6}\right)\leqslant1\)，所以
\[
-3\leqslant f(x)\leqslant1
\]
函数值域为 \([-3,1]\).
\item 当 \(f(x)=-1\) 时，\(\sin\left(\omega x-\frac{\pi}{6}\right)=0\)，所以交点横坐标满足
\[
\omega x-\frac{\pi}{6}=k\pi,
\qquad
x=\frac{k\pi+\pi/6}{\omega},
\qquad k\in\Z
\]
相邻交点间的距离为 \(\frac{\pi}{\omega}\). 由题设 \(\frac{\pi}{\omega}=\frac{\pi}{2}\)，得 \(\omega=2\). 于是
\[
f'(x)=4\cos\left(2x-\frac{\pi}{6}\right)
\]
当 \(-\frac{\pi}{2}+2k\pi<2x-\frac{\pi}{6}<\frac{\pi}{2}+2k\pi\) 时，\(f'(x)>0\)，故函数在
\[
\left(-\frac{\pi}{6}+k\pi,\frac{\pi}{3}+k\pi\right)
\]
上单调递增；当 \(\frac{\pi}{2}+2k\pi<2x-\frac{\pi}{6}<\frac{3\pi}{2}+2k\pi\) 时，\(f'(x)<0\)，故函数在
\[
\left(\frac{\pi}{3}+k\pi,\frac{5\pi}{6}+k\pi\right)
\]
上单调递减，其中 \(k\in\Z\).
\end{enumerate}
\end{solution}

\begin{problem}
已知数列 \(\{a_n\}\)，\(\{b_n\}\) 满足 \(a_1=2\)，\(b_1=1\)，且 \(\begin{cases}
a_n=\frac{3}{4}a_{n-1}+\frac{1}{4}b_{n-1}+1,\\
b_n=\frac{1}{4}a_{n-1}+\frac{3}{4}b_{n-1}+1,
\end{cases}\)（\(n\geqslant2\)）.

\begin{enumerate}
\item 令 \(c_n=a_n+b_n\)，求数列 \(\{c_n\}\) 的通项公式；
\item 求数列 \(\{a_n\}\) 的通项公式及前 \(n\) 项和.
\end{enumerate}
\end{problem}

\begin{solution}
\begin{enumerate}
\item 两个递推式相加，得
\[
c_n=a_n+b_n=a_{n-1}+b_{n-1}+2=c_{n-1}+2
\]
又 \(c_1=a_1+b_1=3\)，所以 \(\{c_n\}\) 是首项为 \(3\)，公差为 \(2\) 的等差数列，从而
\[
c_n=3+2(n-1)=2n+1
\]
\item 令 \(d_n=a_n-b_n\). 两个递推式相减，得
\[
d_n=\frac{1}{2}(a_{n-1}-b_{n-1})=\frac{1}{2}d_{n-1}
\]
因为 \(d_1=a_1-b_1=1\)，所以 \(d_n=\left(\frac{1}{2}\right)^{n-1}\). 由 \(a_n=\frac{c_n+d_n}{2}\)，得
\[
a_n=\frac{2n+1+2^{1-n}}{2}=n+\frac{1}{2}+2^{-n}
\]
因此
\[
\sum_{k=1}^{n}a_k=\sum_{k=1}^{n}\left(k+\frac{1}{2}+2^{-k}\right)=\frac{n(n+1)}{2}+\frac{n}{2}+\left(1-2^{-n}\right)=\frac{n(n+2)}{2}+1-2^{-n}
\]
\end{enumerate}
\end{solution}

\begin{problem}
已知正三角形 \(OAB\) 的三个顶点都在抛物线 \(y^2=2x\) 上，其中 \(O\) 为坐标原点，设圆 \(C\) 是 \(\triangle OAB\) 的外接圆（点 \(C\) 为圆心）.

\begin{enumerate}
\item 求圆 \(C\) 的方程；
\item 设圆 \(M\) 的方程为 \((x-4-7\cos\theta)^2+(y-7\sin\theta)^2=1\)，过圆 \(M\) 上任意一点 \(P\) 分别作圆 \(C\) 的两条切线 \(PE\)，\(PF\)，切点为 \(E\)，\(F\)，求 \(\overrightarrow{CE}\cdot\overrightarrow{CF}\) 的最大值和最小值.
\end{enumerate}
\end{problem}

\begin{solution}
设 \(A=(u^2/2,u)\)，\(B=(v^2/2,v)\)，这是抛物线 \(y^2=2x\) 上点的参数表示.
\begin{enumerate}
\item 因为三角形 \(OAB\) 为正三角形，所以 \(OA=OB\). 于是
\[
\frac{u^4}{4}+u^2=\frac{v^4}{4}+v^2,
\qquad
(u^2-v^2)(u^2+v^2+4)=0
\]
第二个因式不可能为零，故 \(v^2=u^2\). 由于 \(A\neq B\)，有 \(v=-u\). 再由 \(OA=AB\)，得
\[
\frac{u^4}{4}+u^2=(2u)^2,
\qquad u^2=12
\]
所以可以取
\[
A=(6,2\sqrt{3}),\qquad B=(6,-2\sqrt{3})
\]
正三角形的外心也是重心，故圆心
\[
C=\frac{O+A+B}{3}=(4,0)
\]
又 \(CO=4\)，所以圆 \(C\) 的方程为
\[
(x-4)^2+y^2=16
\]
即 \(x^2+y^2-8x=0\).
\item 圆 \(M\) 的圆心记为 \(Q=(4+7\cos\theta,7\sin\theta)\). 由 \(C=(4,0)\) 得 \(CQ=7\)，而圆 \(M\) 的半径为 \(1\)，所以对圆 \(M\) 上的点 \(P\)，有
\[
6\leqslant CP\leqslant8
\]
圆 \(C\) 的半径为 \(4\). 设 \(\angle ECP=\angle FCP=\varphi\)，则在直角三角形 \(CEP\) 中，
\[
\cos\varphi=\frac{CE}{CP}=\frac{4}{CP}
\]
由于 \(\overrightarrow{CE}\) 与 \(\overrightarrow{CF}\) 的夹角为 \(2\varphi\)，所以
\[
\overrightarrow{CE}\cdot\overrightarrow{CF}
=16\cos2\varphi
=16\left(2\cos^2\varphi-1\right)
=\frac{512}{CP^2}-16
\]
该式随 \(CP\) 增大而减小. 因此当 \(CP=6\) 时取得最大值，当 \(CP=8\) 时取得最小值，分别为
\[
\max\left(\overrightarrow{CE}\cdot\overrightarrow{CF}\right)=\frac{512}{36}-16=-\frac{16}{9}
\]
\[
\min\left(\overrightarrow{CE}\cdot\overrightarrow{CF}\right)=\frac{512}{64}-16=-8
\]
\end{enumerate}
\end{solution}

\begin{problem}
已知函数 \(f(x)=x^2-9x^2\cos\alpha+48x\cos\beta+18\sin^2\alpha\)，\(g(x)=f'(x)\)，且对任意的实数 \(t\) 均有 \(g(1+\cos t)\geqslant0\)，\(g(3+\sin t)\leqslant0\).

\begin{enumerate}
\item 求函数 \(f(x)\) 的解析式；
\item 若对任意的 \(m\in[-26,6]\)，恒有 \(f(x)\geqslant x^2-mx-11\)，求 \(x\) 的取值范围.
\end{enumerate}
\end{problem}

\begin{solution}
由题设，
\[
g(x)=2(1-9\cos\alpha)x+48\cos\beta
\]
取 \(t=0\) 代入第一个不等式，取 \(t=\frac{3\pi}{2}\) 代入第二个不等式，得到 \(g(2)\geqslant0\) 且 \(g(2)\leqslant0\)，所以 \(g(2)=0\). 令 \(r=1-9\cos\alpha\)，则
\[
g(x)=2r(x-2)
\]
要使 \(g(x)\geqslant0\) 对所有 \(x\in[0,2]\) 成立，并使 \(g(x)\leqslant0\) 对所有 \(x\in[2,4]\) 成立，只需且必须有 \(r\leqslant0\). 结合 \(g(2)=0\)，可得
\[
\cos\alpha\geqslant\frac{1}{9}
\qquad
48\cos\beta=-4r=4(9\cos\alpha-1)
\]
这些条件不能唯一确定 \(\alpha\)，\(\beta\)，也不能唯一确定 \(f(x)\). 例如：
\[
\cos\alpha=\frac{1}{9},\quad \cos\beta=0
\]
时，\(g(x)=0\)，且
\[
f(x)=18\left(1-\frac{1}{81}\right)=\frac{160}{9}
\]
而
\[
\cos\alpha=\frac{1}{3},\quad \cos\beta=\frac{1}{6}
\]
时，\(g(x)=-4x+8\)，同样满足题设两个符号条件，但
\[
f(x)=-2x^2+8x+16
\]
两种合法结果不同，因此第（1）问按原题条件无唯一解析式.

第（2）问也随所取参数而变，不能得到唯一的 \(x\) 的取值范围. 比如取 \(x=\frac{11}{10}>0\)，对 \(m\in[-26,6]\)，右端 \(x^2-mx-11\) 在 \(m=-26\) 时最大，最大值为
\[
\left(\frac{11}{10}\right)^2+26\cdot\frac{11}{10}-11=\frac{1881}{100}
\]
对于第一组参数，\(f(x)=\frac{160}{9}<\frac{1881}{100}\)，不满足不等式；对于第二组参数，
\[
f\left(\frac{11}{10}\right)=-2\left(\frac{11}{10}\right)^2+8\cdot\frac{11}{10}+16=\frac{2238}{100}>\frac{1881}{100}
\]
满足对所有 \(m\in[-26,6]\) 的不等式. 因此原题第（2）问按现有题面也不存在唯一的 \(x\) 的取值范围，本题保留这一参数不唯一的结论.
\end{solution}
